\documentclass[12pt]{article}
\usepackage[margin=1in]{geometry}
\usepackage{parskip}
\usepackage{longtable}
\usepackage{array}
\usepackage{booktabs}
\usepackage{enumitem}
\usepackage{hyperref}

\hypersetup{
    colorlinks=false,
    hidelinks
}

\title{\textbf{GreenStep Employee Sustainability Challenge}\\
\large Product Requirements Document}
\author{Hadi Shaar, Naisha Srivastav, Oli Gurmessa, and Nathanael Agbemadon}
\date{Spring 2026 \\ Professor Priscilla Jimenez}

\begin{document}

\maketitle

\section{Product Overview}

GreenStep Employee is a web-based sustainability challenge platform that helps
employees log eco-friendly actions, earn points and badges, participate in
teams, and compare progress through leaderboards. The product replaces
spreadsheet-based tracking with a centralized, low-maintenance system that is
more engaging for end users and easier to manage for administrators.

The primary user groups are:
\begin{itemize}[leftmargin=1.5em]
    \item \textbf{Employees}, who join challenges, log actions, and monitor
    their own progress.
    \item \textbf{Administrators}, who configure challenges, review team
    requests, and export participation data.
\end{itemize}

\section{Core Functional Scope}

At the end of the Spring 2026 capstone, the product supports:
\begin{itemize}[leftmargin=1.5em]
    \item secure sign-in and protected dashboard access,
    \item joining challenges through a join code,
    \item joining or leaving a team within a challenge,
    \item submitting team requests for administrator review,
    \item logging sustainability actions across multiple categories,
    \item awarding points, badges, streaks, and leaderboard positions,
    \item viewing challenge-specific and global activity data,
    \item category and challenge filtering for dashboard insights, and
    \item exporting CSV data from the admin workspace.
\end{itemize}

\section{Detailed Use Cases}

\subsection{Use Case 1: Employee Logs a Sustainability Action}

\textbf{Primary Actor:} Employee

\textbf{Goal:} Record a completed sustainability action and receive the
appropriate points inside the correct challenge context.

\textbf{Preconditions:}
\begin{itemize}[leftmargin=1.5em]
    \item The employee is authenticated.
    \item The employee has completed onboarding.
    \item The employee is on the dashboard or log action page.
\end{itemize}

\textbf{Trigger:} The employee decides to log a completed action such as biking
to work, recycling plastic, or turning off unused lights.

\textbf{Main Success Scenario:}
\begin{enumerate}[leftmargin=1.8em]
    \item The employee opens the \texttt{Log Action} experience.
    \item The employee browses by category, recent actions, or popular actions,
    or uses search to find a relevant action.
    \item The employee selects an action from the list.
    \item The system opens the action detail page and shows the base point
    value.
    \item The employee optionally adds a story or image.
    \item The employee submits the action.
    \item The system validates the request and checks for duplicate rapid
    submissions.
    \item The system stores the action in the database, scoped to the selected
    challenge or the global activity context.
    \item The system updates the employee's points, activity history, streak,
    badges, and relevant leaderboard views.
    \item The system returns the employee to a success state where the updated
    action is visible in activity feeds and dashboard summaries.
\end{enumerate}

\textbf{Alternative / Exception Flows:}
\begin{itemize}[leftmargin=1.5em]
    \item If the employee submits the same action too quickly, the system
    rejects the duplicate submission and shows an error.
    \item If the employee is no longer authenticated, the system blocks the
    action and requires sign-in.
    \item If an optional image or note is omitted, the action is still
    accepted with only the applicable points.
\end{itemize}

\textbf{Postconditions:}
\begin{itemize}[leftmargin=1.5em]
    \item A new action record exists in the database.
    \item The employee's activity feed and points reflect the logged action.
    \item Challenge-specific and global views remain consistent with the new
    data.
\end{itemize}

\subsection{Use Case 2: Administrator Reviews a Team Request and Manages Challenges}

\textbf{Primary Actor:} Administrator

\textbf{Goal:} Manage challenge participation by reviewing pending team
requests and configuring challenge data from the admin workspace.

\textbf{Preconditions:}
\begin{itemize}[leftmargin=1.5em]
    \item The administrator is authenticated.
    \item The administrator's Clerk metadata role is set to
    \texttt{admin}.
\end{itemize}

\textbf{Trigger:} The administrator opens the admin workspace to review a team
request or create/export challenge data.

\textbf{Main Success Scenario:}
\begin{enumerate}[leftmargin=1.8em]
    \item The administrator navigates to the admin page.
    \item The system verifies that the current user has admin access.
    \item The system displays pending team requests, recently reviewed
    requests, challenge summaries, and export controls.
    \item The administrator opens a pending team request.
    \item The administrator reviews the requested team name, requester, and
    challenge context.
    \item The administrator chooses to approve or reject the request.
    \item If approved, the system creates the team if needed, adds the user to
    the team, and updates the request status.
    \item If rejected, the system records the rejection status and reviewer.
    \item The administrator may also create a new challenge, generate demo
    challenges, or export activity data to CSV.
    \item The system revalidates the affected dashboard and admin views so the
    updated state is immediately reflected.
\end{enumerate}

\textbf{Alternative / Exception Flows:}
\begin{itemize}[leftmargin=1.5em]
    \item If a non-admin user attempts to access the admin workspace, the
    system blocks access and displays an admin-only message.
    \item If a team request is invalid or has already been processed, the
    system rejects the action and preserves data integrity.
    \item If the administrator attempts to create a duplicate team name within
    the same challenge, the system prevents the duplicate.
\end{itemize}

\textbf{Postconditions:}
\begin{itemize}[leftmargin=1.5em]
    \item The request is stored with an approved or rejected status.
    \item Relevant team, challenge, and dashboard data remain synchronized.
    \item Administrators retain a recent history of reviewed requests.
\end{itemize}

\section{User Stories with Acceptance Criteria}

\subsection{User Story 1: Join a Challenge}

\textbf{Story:} As an employee, I want to join a challenge with a code so that
my actions and team participation are associated with the correct campaign.

\textbf{Acceptance Criteria:}
\begin{itemize}[leftmargin=1.5em]
    \item The user can open a join challenge modal from the dashboard.
    \item The user can enter a challenge code and submit it.
    \item If the code is valid and the challenge is active, the user is added
    to that challenge.
    \item If the user already joined the challenge, the system prevents a
    duplicate join and shows an explanatory error.
    \item If the challenge has ended or the code is invalid, the system
    displays a clear error message.
\end{itemize}

\subsection{User Story 2: View Filtered Progress by Challenge or Global Scope}

\textbf{Story:} As an employee, I want to switch between all activity and a
specific challenge so that I can understand my performance in the right
context.

\textbf{Acceptance Criteria:}
\begin{itemize}[leftmargin=1.5em]
    \item The dashboard provides a scope selector for all activity and
    joined challenges.
    \item When the user switches scope, points, leaderboard data, and activity
    summaries update to match that scope.
    \item Global views aggregate lifetime data across the system.
    \item Challenge views show only challenge-specific points, actions, and
    team context.
    \item The system preserves consistent labels so the user can distinguish
    global and challenge-scoped information.
\end{itemize}

\subsection{User Story 3: Submit a Team Request}

\textbf{Story:} As an employee without a team, I want to request a new team in
my challenge so that an administrator can review and approve it.

\textbf{Acceptance Criteria:}
\begin{itemize}[leftmargin=1.5em]
    \item The user can open a request form from the team area of the dashboard.
    \item The user must already belong to the challenge before submitting a
    request.
    \item The request captures the desired team name and optional message.
    \item The system prevents duplicate pending requests for the same context.
    \item Submitted requests appear in the user's request history and in the
    admin review queue.
\end{itemize}

\subsection{User Story 4: Export Participation Data}

\textbf{Story:} As an administrator, I want to export sustainability activity
data so that I can analyze participation outside the application and share
results with stakeholders.

\textbf{Acceptance Criteria:}
\begin{itemize}[leftmargin=1.5em]
    \item The admin workspace provides CSV export controls.
    \item The administrator can export all activity or challenge-specific
    activity.
    \item The generated CSV includes enough participation data to support
    reporting and review.
    \item Only authorized administrators can access export actions.
    \item The exported data reflects the latest saved records in the system.
\end{itemize}

\section{Notes for Future Teams}

Future teams should treat action logging, challenge membership, team
membership, and admin request review as tightly connected product areas.
Changes in one area often affect leaderboards, dashboard summaries, and export
behavior. Before extending the system, teams should confirm whether a feature
is global, challenge-scoped, or admin-only and keep that distinction explicit
in both the UI and API design.

\end{document}
